\section{Dificultades, Mejoras y Conclusiones}

\subsection{Dificultades Técnicas}
Durante el diseño e implementación se identificaron retos relevantes:
\begin{enumerate}
    \item \textbf{Concurrencia en VirtualBox:} las invocaciones concurrentes a \texttt{VBoxManage} pueden generar condiciones de carrera con \texttt{VBoxSVC}. Se mitigó centralizando y serializando operaciones críticas en el adaptador del hipervisor.
    \item \textbf{Detección dinámica de IP:} la asignación de red en la VM no es inmediata; se requiere sondeo controlado (con espera/reintentos) apoyado en mecanismos de consulta del invitado para obtener una IP utilizable antes de registrar DNS o desplegar.
    \item \textbf{Consistencia transaccional entre subsistemas:} un fallo en DNS, IP o despliegue puede dejar recursos huérfanos. La orquestación incorpora compensaciones (\textit{rollback}) para preservar coherencia entre VM, DNS y contenido.
    \item \textbf{Contraste entre interfaz y estado real:} en consolas operativas es riesgoso mostrar estados \emph{hardcoded}. La interfaz evolucionó hacia indicadores basados en \texttt{/api/salud} y manejo explícito de estados de carga/error, aunque la verificación visual final depende de la estabilización de la UI y del laboratorio.
\end{enumerate}

\subsection{Trabajo Futuro}
\begin{itemize}
    \item Completar evidencias de laboratorio (aprovisionamiento de \texttt{paisajes}/\texttt{negocios}, resolución DNS y capturas) una vez estabilizada la interfaz.
    \item Conectar datos operativos reales a módulos evolutivos (VHosts/eventos) mediante endpoints dedicados o enriquecimiento de respuestas.
    \item Incorporar controles de cuota/recursos por instancia y endurecimiento de seguridad (validación reforzada de \texttt{.zip}, autenticación/autorización si el entorno lo requiere).
    \item Evaluar TLS/HTTPS en el entorno local y estrategias de observabilidad con correlación de eventos por operación.
\end{itemize}

\subsection{Conclusiones}
Se diseñó e implementó una base funcional de HTTPaaS que integra virtualización, DNS dinámico y orquestación en Go para automatizar el ciclo de vida de servicios HTTP en un entorno local. La arquitectura hexagonal permitió desacoplar el dominio de la infraestructura, favoreciendo pruebas automatizadas y evolución controlada. El contrato OpenAPI y la interfaz web modular proporcionan un punto de acceso coherente para el aprovisionamiento mediante archivos \texttt{.zip}, mientras que la incorporación de TDD, pruebas de contrato y consideraciones de accesibilidad/verdad operativa fortalecen la mantenibilidad. Queda pendiente la consolidación final de evidencias en el laboratorio y la estabilización visual de la interfaz para cerrar la entrega documental.
